\documentclass[11pt]{article}

\usepackage[utf8]{inputenc}
\usepackage[T1]{fontenc}
\usepackage{amsmath, amssymb, amsthm}
\usepackage{hyperref}
\usepackage[margin=1in]{geometry}

\newtheorem{theorem}{Theorem}[section]
\newtheorem{lemma}[theorem]{Lemma}
\newtheorem{proposition}[theorem]{Proposition}
\newtheorem{corollary}[theorem]{Corollary}
\newtheorem{definition}[theorem]{Definition}
\newtheorem{assumption}[theorem]{Assumption}
\newtheorem{remark}[theorem]{Remark}

\newcommand{\E}{\mathbb{E}}
\newcommand{\Prob}{\mathbb{P}}
\newcommand{\R}{\mathbb{R}}
\newcommand{\cH}{\mathcal{H}}
\newcommand{\cU}{\mathcal{U}}
\newcommand{\cV}{\mathcal{V}}

\title{Main Theorem: Asymptotic Consistency of TVD-MI Judge Ensembles\\
\large Natural-Language Proof for Lean 4 Formalization}

\author{Brando Miranda}
\date{\today}

\begin{document}
\maketitle

\begin{abstract}
This document contains the full natural-language statements and proofs of the
main theorems for the TVD-MI judge ensemble consistency paper. It is intended
as the source for Lean~4 formalization in this experiment folder.
\end{abstract}

\tableofcontents

% ============================================================================
\section{Setup and Notation}
\label{sec:setup}

Let $c$ denote the candidate system or output family being evaluated.

Let $(U, V)$ be the two objects whose dependence defines the reference-free
signal. In the simplest prompt--response form, $U = \text{prompt}$ and
$V = \text{response}$. In the Robertson--Koyejo style peer-prediction form,
$(U, V)$ can be two reports tied to the same latent source.

\begin{definition}[Latent TVD-MI target]
\label{def:tvd-mi-target}
For a candidate $c$, define
\begin{equation}
\tau(c)
:= \mathrm{TV}\!\big(P^{(c)}_{UV},\, P^{(c)}_{U}\otimes P^{(c)}_{V}\big)
= \sup_{h \in \cH} \Delta_c(h),
\end{equation}
where $\cH = \{h : \cU \times \cV \to [0,1]\}$ and
\begin{equation}
\Delta_c(h)
:= \E_{P^{(c)}_{UV}}[h(U,V)]
- \E_{P^{(c)}_{U}\otimes P^{(c)}_{V}}[h(U,V)].
\end{equation}
\end{definition}

\begin{definition}[Gold metric and validity link]
\label{def:gold-metric}
Let $g(c)$ be the inaccessible gold metric. We say there is a
\emph{monotone validity link} if
\begin{equation}
\tau(c_1) > \tau(c_2) \iff g(c_1) > g(c_2).
\end{equation}
More generally, $g(c) = \phi(\tau(c))$ for some known or learnable
monotone $\phi$.
\end{definition}

\begin{remark}
The key architectural point: first prove consistency to $\tau(c)$, then
add the validity link from $\tau(c)$ to $g(c)$. This separation is the
heart of the proposal.
\end{remark}

% ============================================================================
\section{Theorem 1: Random-Witness Consistency (Coverage Regime)}
\label{sec:coverage}

\begin{assumption}[Random witness model]
\label{asm:random-witness}
Each judge $j \in \{1, \ldots, K\}$ is a random critic
$h_j \sim \mu$ drawn i.i.d.\ from a judge distribution $\mu$ over
witness functions $h : \cU \times \cV \to [0,1]$.
\end{assumption}

\begin{definition}[Best-of-$K$ ensemble estimator]
\label{def:best-of-K}
\begin{equation}
\hat\tau^{\max}_K(c) := \max_{1 \le j \le K} \Delta_c(h_j).
\end{equation}
\end{definition}

\begin{definition}[Near-optimal witness mass]
\label{def:near-optimal-mass}
For any $\varepsilon > 0$,
\begin{equation}
p_\varepsilon(c) := \mu\big\{h : \Delta_c(h) \ge \tau(c) - \varepsilon\big\}.
\end{equation}
\end{definition}

\begin{theorem}[Random-witness consistency]
\label{thm:coverage}
Under Assumption~\ref{asm:random-witness}, for any $\varepsilon > 0$,
\begin{equation}
\Prob\!\big(\tau(c) - \hat\tau^{\max}_K(c) > \varepsilon\big)
= (1 - p_\varepsilon(c))^K
\le e^{-K \, p_\varepsilon(c)}.
\end{equation}
\end{theorem}

\begin{proof}
The event $\{\tau(c) - \hat\tau^{\max}_K(c) > \varepsilon\}$ is equivalent
to $\{\max_{j \le K} \Delta_c(h_j) < \tau(c) - \varepsilon\}$, which is
equivalent to $\{\Delta_c(h_j) < \tau(c) - \varepsilon \text{ for all }
j = 1, \ldots, K\}$.

Since $h_1, \ldots, h_K$ are i.i.d.\ from $\mu$,
\begin{align}
\Prob\!\big(\tau(c) - \hat\tau^{\max}_K(c) > \varepsilon\big)
&= \prod_{j=1}^K \Prob\!\big(\Delta_c(h_j) < \tau(c) - \varepsilon\big) \\
&= \prod_{j=1}^K \big(1 - p_\varepsilon(c)\big) \\
&= \big(1 - p_\varepsilon(c)\big)^K.
\end{align}

Using $1 - x \le e^{-x}$ for $x \in [0,1]$, we obtain
\begin{equation}
\big(1 - p_\varepsilon(c)\big)^K \le e^{-K \, p_\varepsilon(c)}.
\end{equation}
\end{proof}

\begin{corollary}[Scaling law under polynomial witness mass]
\label{cor:scaling-law-coverage}
If the near-optimal mass scales as $p_\varepsilon(c) \asymp \varepsilon^\beta$
for some $\beta > 0$, then
\begin{equation}
\tau(c) - \hat\tau^{\max}_K(c) = O_p(K^{-1/\beta}).
\end{equation}
\end{corollary}

\begin{proof}
Set $\varepsilon = K^{-1/\beta}$. Then $p_\varepsilon(c) \asymp K^{-1}$,
and $e^{-K \cdot K^{-1}} = e^{-1}$, which is a constant. So the probability
that the gap exceeds $K^{-1/\beta}$ is bounded by a constant, establishing
the $O_p$ rate.
\end{proof}

\begin{remark}
The exponent $\beta$ encodes the geometry of the judge distribution near
optimal critics. The judge-scaling law is \emph{not} universal; it depends
on diversity.
\end{remark}

% ============================================================================
\section{Theorem 2: Consensus-to-Optimal-Critic (Consensus Regime)}
\label{sec:consensus}

\begin{definition}[Bayes-optimal binary critic]
\label{def:bayes-critic}
Let $h^\star$ be the Bayes-optimal binary critic achieving the total
variation distance, i.e.,
\begin{equation}
\Delta_c(h^\star) = \tau(c).
\end{equation}
\end{definition}

\begin{assumption}[Positive competence]
\label{asm:competence}
For each evaluation pair $z = (u, v)$, judge $j$ outputs a binary vote
$b_j(z) \in \{0, 1\}$ with
\begin{equation}
\Prob\!\big(b_j(z) = h^\star(z) \mid z\big) \ge \tfrac{1}{2} + \gamma(z),
\end{equation}
where $\gamma(z) > 0$ is a pointwise competence margin. Votes are
conditionally independent across judges given $z$.
\end{assumption}

\begin{definition}[Majority vote]
\label{def:majority-vote}
\begin{equation}
b_K(z) := \mathbf{1}\!\left[\sum_{j=1}^K b_j(z) > K/2\right].
\end{equation}
\end{definition}

\begin{theorem}[Consensus convergence]
\label{thm:consensus}
Under Assumption~\ref{asm:competence}, let $M$ denote the mixture
distribution over evaluation pairs. Then
\begin{equation}
0 \le \tau(c) - \Delta_c(b_K) \le 2 \, \E_{Z \sim M}\!\left[
e^{-2K\gamma(Z)^2}\right].
\end{equation}
If $\gamma(z) \ge \gamma_0 > 0$ almost surely, then
\begin{equation}
\tau(c) - \Delta_c(b_K) \le 2 e^{-2K\gamma_0^2}.
\end{equation}
\end{theorem}

\begin{proof}
Fix $z = (u,v)$. Let $X_j = \mathbf{1}[b_j(z) = h^\star(z)]$, so each
$X_j$ is Bernoulli with parameter $q(z) \ge \frac{1}{2} + \gamma(z)$.
The variables $X_1, \ldots, X_K$ are independent given $z$.

The majority vote $b_K(z) \neq h^\star(z)$ iff fewer than $K/2$ judges
agree with $h^\star(z)$, i.e., $\sum_j X_j \le K/2$.

By Hoeffding's inequality applied to the sum $\sum_j X_j$ with
$\E[\sum_j X_j \mid z] = K q(z) \ge K(\frac{1}{2} + \gamma(z))$:
\begin{equation}
\Prob\!\big(b_K(z) \neq h^\star(z) \mid z\big)
\le \Prob\!\left(\sum_j X_j \le K/2 \mid z\right)
\le e^{-2K\gamma(z)^2}.
\end{equation}

Now, since $h^\star(z) \in \{0,1\}$ and $b_K(z) \in \{0,1\}$:
\begin{align}
\Delta_c(h^\star) - \Delta_c(b_K)
&= \E_M\!\big[h^\star(Z) - b_K(Z)\big]_{\text{joint}}
  - \E_M\!\big[h^\star(Z) - b_K(Z)\big]_{\text{product}} \\
&\le 2 \, \E_M\!\big[|h^\star(Z) - b_K(Z)|\big] \\
&= 2 \, \E_M\!\big[\Prob(b_K(Z) \neq h^\star(Z) \mid Z)\big] \\
&\le 2 \, \E_M\!\big[e^{-2K\gamma(Z)^2}\big].
\end{align}

Under the uniform margin $\gamma(z) \ge \gamma_0$:
\begin{equation}
2 \, \E_M\!\big[e^{-2K\gamma(Z)^2}\big]
\le 2 \, e^{-2K\gamma_0^2}.
\end{equation}
\end{proof}

\begin{remark}
Under positive competence and independence, majority-vote judges converge
\emph{exponentially fast} to the optimal TVD critic. This is the cleanest
formalization of the asymptotic-consensus hypothesis.
\end{remark}

% ============================================================================
\section{Theorem 3: Two-Axis Error Decomposition (Judges $\times$ Samples)}
\label{sec:two-axis}

\begin{assumption}[Finite-sample estimation]
\label{asm:finite-sample}
Each judge $j$'s critic value $\Delta_c(h_j)$ is estimated empirically from
$n$ i.i.d.\ positive/negative examples, giving $\widehat\Delta_c(h_j)$.
\end{assumption}

\begin{theorem}[Two-axis decomposition]
\label{thm:two-axis}
Under Assumptions~\ref{asm:random-witness} and~\ref{asm:finite-sample},
for any $\varepsilon > 0$ and $\eta > 0$:
\begin{equation}
\Prob\!\Big(\tau(c) - \max_{j \le K} \widehat\Delta_c(h_j) > \varepsilon + \eta\Big)
\le e^{-K \, p_\varepsilon(c)} + 2K \, e^{-2n\eta^2}.
\end{equation}
\end{theorem}

\begin{proof}
Decompose the event via a union bound. Define the ``good'' event
$G = \{|\widehat\Delta_c(h_j) - \Delta_c(h_j)| \le \eta \text{ for all }
j = 1, \ldots, K\}$.

\textbf{Step 1: Concentration per judge.}
Since $h_j$ maps to $[0,1]$ and $\Delta_c(h_j)$ is an expectation of a
$[0,1]$-bounded random variable, Hoeffding's inequality gives
\begin{equation}
\Prob\!\big(|\widehat\Delta_c(h_j) - \Delta_c(h_j)| > \eta\big)
\le 2 e^{-2n\eta^2}.
\end{equation}

By a union bound over $K$ judges:
\begin{equation}
\Prob(G^c) \le 2K \, e^{-2n\eta^2}.
\end{equation}

\textbf{Step 2: Judge approximation conditional on $G$.}
On event $G$, for every $j$:
$\widehat\Delta_c(h_j) \ge \Delta_c(h_j) - \eta$.

Therefore
\begin{equation}
\max_{j \le K} \widehat\Delta_c(h_j)
\ge \max_{j \le K} \Delta_c(h_j) - \eta
= \hat\tau^{\max}_K(c) - \eta.
\end{equation}

So on $G$:
\begin{equation}
\tau(c) - \max_{j \le K} \widehat\Delta_c(h_j)
\le \tau(c) - \hat\tau^{\max}_K(c) + \eta.
\end{equation}

If this exceeds $\varepsilon + \eta$, then
$\tau(c) - \hat\tau^{\max}_K(c) > \varepsilon$, which by
Theorem~\ref{thm:coverage} has probability at most
$(1 - p_\varepsilon(c))^K \le e^{-K \, p_\varepsilon(c)}$.

\textbf{Step 3: Combine.}
\begin{align}
&\Prob\!\Big(\tau(c) - \max_{j \le K} \widehat\Delta_c(h_j) > \varepsilon + \eta\Big) \\
&\quad\le \Prob\!\Big(\tau(c) - \hat\tau^{\max}_K(c) > \varepsilon\Big) + \Prob(G^c) \\
&\quad\le e^{-K \, p_\varepsilon(c)} + 2K \, e^{-2n\eta^2}.
\end{align}
\end{proof}

\begin{remark}
This cleanly separates:
\begin{itemize}
\item \textbf{Judge approximation error}: controlled by $K$ and the
  witness mass $p_\varepsilon(c)$;
\item \textbf{Finite-sample estimation error}: controlled by $n$ per judge.
\end{itemize}
\end{remark}

% ============================================================================
\section{Judge Dependence and Diversity}
\label{sec:dependence}

\begin{definition}[Equicorrelated judge scores]
\label{def:equicorrelated}
Let $S_j$ be scalar judge scores for a fixed candidate $c$, with
\begin{equation}
\mathrm{Var}(S_j \mid c) = \sigma^2,
\qquad
\mathrm{Corr}(S_i, S_j \mid c) = \rho_K(c) \quad \text{for } i \neq j.
\end{equation}
\end{definition}

\begin{proposition}[Variance of the ensemble mean]
\label{prop:variance-floor}
Let $\bar S_K := \frac{1}{K} \sum_{j=1}^K S_j$. Then
\begin{equation}
\mathrm{Var}(\bar S_K \mid c)
= \sigma^2 \left(\frac{1 - \rho_K(c)}{K} + \rho_K(c)\right).
\end{equation}
\end{proposition}

\begin{proof}
\begin{align}
\mathrm{Var}(\bar S_K \mid c)
&= \frac{1}{K^2} \sum_{i,j} \mathrm{Cov}(S_i, S_j \mid c) \\
&= \frac{1}{K^2} \Big[K \sigma^2 + K(K-1) \rho_K(c) \sigma^2\Big] \\
&= \frac{\sigma^2}{K} \Big[1 + (K-1)\rho_K(c)\Big] \\
&= \sigma^2 \left(\frac{1 - \rho_K(c)}{K} + \rho_K(c)\right).
\end{align}
\end{proof}

\begin{corollary}[Scaling regimes]
\label{cor:scaling-regimes}
\begin{itemize}
\item If $\rho_K(c) = 0$: $\mathrm{RMSE} \sim K^{-1/2}$.
\item If $\rho_K(c) = \rho_0 > 0$ (constant): irreducible error floor
  $\sigma\sqrt{\rho_0}$.
\item If $\rho_K(c) \sim K^{-\beta}$ for $\beta > 0$: $\mathrm{RMSE} \sim
  K^{-\alpha}$ with $\alpha = \frac{1}{2}\min(1, \beta)$.
\end{itemize}
\end{corollary}

\begin{theorem}[Impossibility under persistent dependence]
\label{thm:impossibility-dependence}
If residual pairwise dependence does not vanish, i.e.,
$\liminf_{K \to \infty} \rho_K(c) > 0$, then averaging-based aggregation
cannot be asymptotically consistent:
\begin{equation}
\liminf_{K \to \infty} \mathrm{Var}(\bar S_K \mid c) > 0.
\end{equation}
\end{theorem}

\begin{proof}
Immediate from Proposition~\ref{prop:variance-floor}: if
$\rho_K(c) \ge \rho_0 > 0$ for all $K$, then
$\mathrm{Var}(\bar S_K \mid c) \ge \sigma^2 \rho_0 > 0$.
\end{proof}

% ============================================================================
\section{Adversarial Robustness}
\label{sec:adversarial}

\begin{assumption}[Contamination model]
\label{asm:contamination}
A fraction $(1-\delta)$ of judges are honest with competence margin
$\gamma > 0$. A fraction $\delta$ are adversarial.
\end{assumption}

\begin{theorem}[Adversarial breakdown threshold]
\label{thm:adversarial}
For binary votes where honest judges are correct with probability
$\frac{1}{2} + \gamma$ and adversaries are always wrong:
\begin{itemize}
\item The ensemble is asymptotically correct iff
\begin{equation}
(1-\delta)\Big(\tfrac{1}{2} + \gamma\Big) > \tfrac{1}{2}
\quad\Longleftrightarrow\quad
\delta < \frac{2\gamma}{1 + 2\gamma}.
\end{equation}
\item Under the milder sign-flip adversary model, the threshold becomes
$\delta < \gamma / (1 + \gamma)$.
\end{itemize}
\end{theorem}

\begin{proof}
The expected fraction of correct votes is
\begin{equation}
(1-\delta)\Big(\tfrac{1}{2} + \gamma\Big) + \delta \cdot 0
= (1-\delta)\Big(\tfrac{1}{2} + \gamma\Big).
\end{equation}

Majority vote succeeds asymptotically iff this exceeds $\frac{1}{2}$:
\begin{align}
(1-\delta)\Big(\tfrac{1}{2} + \gamma\Big) &> \tfrac{1}{2} \\
\tfrac{1}{2} + \gamma - \delta\Big(\tfrac{1}{2} + \gamma\Big) &> \tfrac{1}{2} \\
\gamma &> \delta\Big(\tfrac{1}{2} + \gamma\Big) \\
\delta &< \frac{\gamma}{\frac{1}{2} + \gamma} = \frac{2\gamma}{1 + 2\gamma}.
\end{align}

For sign-flip adversaries (correct with probability $\frac{1}{2} - \gamma$):
\begin{align}
(1-\delta)\Big(\tfrac{1}{2} + \gamma\Big) + \delta\Big(\tfrac{1}{2} - \gamma\Big) &> \tfrac{1}{2} \\
\tfrac{1}{2} + \gamma - 2\delta\gamma &> \tfrac{1}{2} \\
\delta &< \frac{1}{2}.
\end{align}

Wait---let me redo this more carefully. For sign-flip adversaries who are
correct with probability $\frac{1}{2} - \gamma$:
\begin{align}
(1-\delta)\Big(\tfrac{1}{2} + \gamma\Big) + \delta\Big(\tfrac{1}{2} - \gamma\Big) &> \tfrac{1}{2} \\
\tfrac{1}{2} + \gamma(1-2\delta) &> \tfrac{1}{2} \\
\gamma(1-2\delta) &> 0 \\
\delta &< \tfrac{1}{2}.
\end{align}

So the sign-flip threshold is $\delta < 1/2$, which is the standard
Byzantine majority threshold.
\end{proof}

\begin{proposition}[Robust aggregation error]
\label{prop:robust-aggregation}
For bounded scalar scores in $[0,1]$, any robust estimator inherits an
unavoidable $O(\delta)$ bias term. A trimmed mean or median-of-means gives
\begin{equation}
\mathrm{Error} = O_p(K^{-1/2} + \delta).
\end{equation}
\end{proposition}

\begin{remark}[Optimistic vs.\ pessimistic regimes]
\begin{itemize}
\item \textbf{Optimistic}: $\rho \to 0$ and $\delta$ below breakdown
  $\Longrightarrow$ consistency.
\item \textbf{Pessimistic}: correlation floor or too much contamination
  $\Longrightarrow$ irreducible error.
\end{itemize}
\end{remark}

% ============================================================================
\section{Gold-Metric Consistency}
\label{sec:gold}

\begin{theorem}[Rank consistency with gold metric]
\label{thm:gold-consistency}
Suppose a monotone validity link holds (Definition~\ref{def:gold-metric}).
If the ensemble estimator $\hat\tau_K(c)$ is consistent for $\tau(c)$
uniformly over a finite set of candidates $\{c_1, \ldots, c_L\}$, then
the induced ranking converges to the gold ranking.
\end{theorem}

\begin{proof}
By uniform consistency, for any $\varepsilon > 0$,
\begin{equation}
\Prob\!\Big(\max_{l \le L} |\hat\tau_K(c_l) - \tau(c_l)| > \varepsilon\Big) \to 0
\quad\text{as } K \to \infty.
\end{equation}

Let $\Delta_{\min} = \min_{l \neq l'} |\tau(c_l) - \tau(c_{l'})|$ be the
minimum gap between distinct candidates (assumed positive). For
$\varepsilon < \Delta_{\min}/2$, the event
$\{\max_l |\hat\tau_K(c_l) - \tau(c_l)| \le \varepsilon\}$ implies
the ranking induced by $\hat\tau_K$ agrees with the ranking induced by
$\tau$. By the monotone link, this agrees with the gold ranking.
\end{proof}

% ============================================================================
\section{Summary of Main Results}
\label{sec:summary}

The strongest first-paper result is the three-part theorem:

\begin{enumerate}
\item \textbf{Coverage regime} (Theorem~\ref{thm:coverage}):
\[
\Prob\!\Big(\tau(c) - \max_{j \le K} \widehat\Delta_c(h_j)
> \varepsilon + \eta\Big)
\le e^{-K \, p_\varepsilon(c)} + 2K \, e^{-2n\eta^2}.
\]

\item \textbf{Diversity as covariance} (Proposition~\ref{prop:variance-floor}):
\[
\mathrm{Var}(\bar S_K \mid c)
= \sigma^2\!\left(\frac{1 - \rho_K(c)}{K} + \rho_K(c)\right).
\]

\item \textbf{Impossibility modes}:
\begin{itemize}
\item Persistent dependence $\Rightarrow$ irreducible variance floor.
\item Contamination $\delta \ge 2\gamma/(1+2\gamma)$ $\Rightarrow$ breakdown.
\end{itemize}
\end{enumerate}

This gives: a TVD-specific judge-scaling theorem, a diversity-as-covariance
definition, a clear impossibility mode, and an empirical recipe.

\end{document}
